\section{Fixed-block coercivity in the actual vacuum}
\label{vac:section}

This section proves a lower bound for the energy of fluctuations in a fixed
finite set of links, conditional on all exterior links, in the actual
finite-volume ground state.  The constant is independent of the number of
exterior links.  Its dependence on the number of selected links is explicit
and does not give a uniform bound for an extensive elimination layer.
We use a Hamiltonian of Kogut--Susskind type~\cite{KS} with the normalization
specified below.  The argument is presented in full; the associated research
record is available in~\cite{Archive}.

\subsection{Finite-volume Hamiltonian and its ground-state form}
\label{vac:setup}

Let $E$ be a finite set of oriented link representatives of a finite spatial
graph, let $\mathcal P$ be a finite collection of plaquette terms, and fix
$N\geq2$, $x>0$, $a>0$, and $b\geq0$.  Put
$M=\SU(N)^E$ and equip $M$ with normalized product Haar measure $\mathrm dU$.
For Hermitian generators satisfying
$\operatorname{Tr}(t_\alpha t_\beta)=\delta_{\alpha\beta}/2$, choose the
bi-invariant metric for which the corresponding left-invariant vector fields
$X_\alpha$ are orthonormal.  On link $\ell$, write
$C_\ell=-\sum_\alpha X_{\ell,\alpha}^2=-\Delta_\ell$.  Consider
\begin{equation}
 H=\frac{x}{2}\sum_{\ell\in E}C_\ell+V,
 \qquad
 V(U)=\frac b x\sum_{p\in\mathcal P}
       \left(1-\frac{\operatorname{Re}\operatorname{Tr}U_p}{N}\right).
 \label{vac:H}
\end{equation}
Here $U_p$ is the ordered plaquette holonomy.  Each summand in parentheses
lies in $[0,2]$.  The operator $H$ is the self-adjoint bounded-potential
perturbation of the product Laplacian, with operator domain $H^2(M)$ and
form domain $H^1(M)$.  The parameter $a$ enters the physical energy through
division by $a$; it has not been absorbed into $H$.

The vertex gauge group acts by
$U_\ell\mapsto g_{s(\ell)}U_\ell g_{t(\ell)}^{-1}$.  We assume the
plaquette terms and boundary conditions preserve this action.  Physical
functions are invariant under this group.  No gauge fixing is used.

\begin{lemma}[Actual vacuum and ground-state transform]
\label{vac:groundstate}
The lowest eigenvalue $E_0$ of $H$ is simple.  Its normalized eigenfunction
$\psi$ can be chosen smooth, strictly positive, and gauge invariant.  Set
$\mathrm d\mu=\psi^2\,\mathrm dU$.  Multiplication by $\psi$ is a unitary map
from $L^2(\mu)$ to $L^2(\mathrm dU)$, and the nonnegative self-adjoint operator
\begin{equation}
 K=\psi^{-1}(H-E_0)\psi/a
 \label{vac:K}
\end{equation}
has closed quadratic form
\begin{equation}
 k[f]=\frac{x}{2a}\int_M
          \sum_{\ell,\alpha}|X_{\ell,\alpha}f|^2\,\mathrm d\mu,
 \qquad f\in H^1(M).
 \label{vac:groundform}
\end{equation}
Its kernel is the space of constant functions.  These statements also hold
after restriction to the gauge-invariant subspaces.
\end{lemma}

\begin{proof}
The compact manifold $M$ is connected, whether or not the underlying graph
is connected.  Its heat kernel is strictly positive at every positive time.
The Feynman--Kac weight for the bounded real potential $V$ is strictly
positive, so $\exp(-tH)$ is positivity improving.  Compact resolvent gives a
lowest eigenfunction.  The absolute-value inequality for the Dirichlet form
allows a nonnegative minimizer, which is strictly positive by the strong
maximum principle.  Positivity improvement implies simplicity: a second
real ground eigenfunction orthogonal to the positive one would change sign,
contradicting the strict positivity statement.  Elliptic regularity gives
smoothness.  Gauge transformations commute with $H$ and preserve positivity
and norm; uniqueness therefore makes $\psi$ gauge invariant.

For smooth complex $f$, integration by parts and $H\psi=E_0\psi$ give
\begin{equation}
 \langle\psi f,(H-E_0)\psi f\rangle_{L^2(\mathrm dU)}
   =\frac{x}{2}\int_M
       \sum_{\ell,\alpha}|X_{\ell,\alpha}f|^2\psi^2\,\mathrm dU.
 \label{vac:transformidentity}
\end{equation}
Indeed, the terms containing $|f|^2|\nabla\psi|^2$ and the mixed derivatives
combine, after integration by parts, with $V|\psi f|^2$ to give
$E_0|\psi f|^2$.  On this fixed finite graph, $\psi$ is smooth and bounded
away from zero.  Multiplication by $\psi$ and its inverse preserve $H^1(M)$,
so the identity extends by form closure.  Unitarity follows directly from
$\mathrm d\mu=\psi^2\mathrm dU$.  The zero-energy functions have zero gradient
and are constant because $M$ is connected.  Multiplication by $\psi$
intertwines the gauge actions, which proves the final assertion.
\end{proof}

In particular, all energy estimates below subtract $E_0$, the actual
interacting ground energy.  No bound on $E_0$ uniform in volume is assumed.
The equivalence of weighted and unweighted Sobolev norms just used is needed
only for each fixed graph; its constants need not be uniform in volume.

\subsection{Conditional projection and finite-block theorem}
\label{vac:block}

Fix a nonempty subset $B\subseteq E$ with $m=|B|$, and write a configuration
as $(u,y)\in\SU(N)^B\times\SU(N)^{E\setminus B}$.  Let $p_B$ be the number
of plaquette terms that involve at least one link in $B$, counting each term
once.  Define
\begin{equation}
 Z(y)=\int\psi(v,y)^2\,\mathrm dv,
 \qquad
 \mathrm d\mu_y(u)=\frac{\psi(u,y)^2}{Z(y)}\,\mathrm du,
 \qquad
 Pf(u,y)=\int f(v,y)\,\mathrm d\mu_y(v),
 \label{vac:conditional}
\end{equation}
where $\mathrm du$ and $\mathrm dv$ are normalized product Haar measures.
Thus $P=\mathbb E_\mu[\,\cdot\mid y]$ is an orthogonal projection in
$L^2(\mu)$ and $P1=1$.  Put $Q=I-P$.

The restriction of the form $k$ to
$H^1(M)\cap\operatorname{Ran}Q$ is closed and densely defined in
$\operatorname{Ran}Q$; the density assertion is verified in the proof below.
Let $D_B$ be its associated self-adjoint operator on that subspace.  Thus
$D_B$ denotes a form compression.  It does not require the unbounded
operator product $QKQ$ to be defined on an assumed common operator domain.

Write
\begin{equation}
 C_F=\frac{N^2-1}{2N},\qquad n=N^2-1,
 \qquad
 \kappa_N(s)=
 \frac{\max_{u,v\in\SU(N)}k_s(u,v)}
      {\min_{u,v\in\SU(N)}k_s(u,v)},\quad s>0,
 \label{vac:constants}
\end{equation}
where $k_s$ is the one-link heat kernel of $\exp(-sC)=\exp(s\Delta)$
relative to normalized Haar measure.  Compactness and strict positivity make
$\kappa_N(s)$ finite.

\begin{theorem}[Exterior-volume-uniform finite-block coercivity]
\label{vac:finiteblock}
For every $s>0$,
\begin{equation}
 D_B\ \geq\ d_B(s)I,
 \qquad
 d_B(s)=\frac{xC_F}{2a}\,
       \kappa_N(s)^{-2m}
       \exp\!\left(-\frac{8sbp_B}{x^2}\right).
 \label{vac:heatbound}
\end{equation}
The bound is independent of the number of links outside $B$ when
$m,p_B,N,x,b,a$ are fixed.  The projection $P$ commutes with the gauge
action, so the same bound holds on the gauge-invariant part of
$\operatorname{Ran}Q$, if that part is nonzero.  If $bp_B=0$, the sharper
bound
\begin{equation}
 D_B\geq\frac{xC_F}{2a}I
 \label{vac:freebound}
\end{equation}
holds.
\end{theorem}

\begin{proof}
Split the potential into terms that avoid and meet $B$:
\begin{equation}
 V(u,y)=V_{\mathrm{out}}(y)+V_{\mathrm{hit}}(u,y),
 \qquad
 0\leq V_{\mathrm{hit}}\leq W_B:=\frac{2bp_B}{x}.
 \label{vac:split}
\end{equation}
Let $(U_t,Y_t)$ be independent link Brownian motions, with generator
$(x/2)\Delta$ on each link.  At any time $t>0$, the eigenfunction identity
and Feynman--Kac formula give
\begin{equation}
 e^{-tE_0}\psi(u,y)=
 \mathbb E_{u,y}\!\left[
 e^{-\int_0^t V(U_r,Y_r)\,\mathrm dr}\psi(U_t,Y_t)\right].
 \label{vac:fk}
\end{equation}
Remove only $V_{\mathrm{hit}}$ from this expectation and call the result
$F_t(u,y)$.  Positivity gives
\begin{equation}
 e^{-tW_B}F_t(u,y)\leq e^{-tE_0}\psi(u,y)\leq F_t(u,y).
 \label{vac:fksandwich}
\end{equation}
Put $s=xt/2$.  Independence of the link motions and the absence of $u$ from
$V_{\mathrm{out}}$ imply
\begin{align}
 F_t(u,y)&=\int k_s^B(u,v)A_t(y,v)\,\mathrm dv,
 \label{vac:factorization}\\
 A_t(y,v)&=\mathbb E_y\!\left[
       e^{-\int_0^t V_{\mathrm{out}}(Y_r)\,\mathrm dr}
       \psi(v,Y_t)\right]>0,
 \notag
\end{align}
where $k_s^B$ is the product of the $m$ one-link kernels.  For a fixed $y$,
the positive factor $A_t(y,v)$ is the same for every initial block
configuration $u$.  Consequently, for every $u,u',y$,
\begin{equation}
 \frac{\psi(u,y)}{\psi(u',y)}
 \leq e^{tW_B}\frac{F_t(u,y)}{F_t(u',y)}
 \leq e^{tW_B}\kappa_N(s)^m=:R_B(s).
 \label{vac:oscillation}
\end{equation}
The factor $e^{-tE_0}$ cancels in this same-time ratio.  The exterior
potential remains inside the common positive expectation in
\eqref{vac:factorization}; it has not been replaced by a bound depending on
the number of exterior plaquettes.  This is the source of the asserted
exterior-volume independence.

We next recall the product Haar Poincar\'e inequality in the chosen metric:
\begin{equation}
 \operatorname{Var}_{\mathrm du}(f)
 \leq\frac1{C_F}\int\sum_{\ell\in B,\alpha}
                 |X_{\ell,\alpha}f|^2\,\mathrm du.
 \label{vac:haar}
\end{equation}
For completeness, Peter--Weyl decomposition reduces this statement to the
smallest nonzero Casimir eigenvalue of $\SU(N)$.  With long roots of squared
length $2$, the Casimir of highest weight $\lambda$ is
$(\lambda,\lambda+2\rho)/2$.  The fundamental weight $\omega_j$ has Casimir
$j(N-j)(N+1)/(2N)$.  All pairings of fundamental weights are nonnegative,
so every nonzero dominant integral weight has Casimir at least the minimum
of these fundamental values, namely $C_F$.  On a product, the eigenvalues
add; a nonconstant component therefore has eigenvalue at least $C_F$.
This proves~\eqref{vac:haar}, including for complex functions by applying
the real inequality to their real and imaginary parts.

For any positive probability density $w$ relative to product Haar with
$\max w/\min w\leq R^2$, the minimizing-constant definition of variance and
\eqref{vac:haar} give
\begin{align}
 \operatorname{Var}_{w\mathrm du}(f)
 &\leq (\max w)\operatorname{Var}_{\mathrm du}(f)\notag\\
 &\leq\frac{\max w}{C_F}\int|\nabla_Bf|^2\,\mathrm du
 \leq\frac{R^2}{C_F}\int|\nabla_Bf|^2w\,\mathrm du,
 \label{vac:densitycomparison}
\end{align}
where $|\nabla_Bf|^2=\sum_{\ell\in B,\alpha}|X_{\ell,\alpha}f|^2$.
No Haar-mean-zero condition is imposed on $f$ in this comparison.
By~\eqref{vac:oscillation}, the density of $\mu_y$ in
\eqref{vac:conditional} has maximum-to-minimum ratio at most $R_B(s)^2$.
Applying~\eqref{vac:densitycomparison} conditionally and integrating over
the exterior marginal yields
\begin{equation}
 \|Qf\|_{L^2(\mu)}^2
 \leq\frac{R_B(s)^2}{C_F}
        \int_M|\nabla_Bf|^2\,\mathrm d\mu.
 \label{vac:conditionalpoincare}
\end{equation}
If $f\in\operatorname{Ran}Q$, the exact ground-state form therefore gives
\begin{equation}
 k[f]\geq\frac{x}{2a}\int_M|\nabla_Bf|^2\,\mathrm d\mu
       \geq\frac{xC_F}{2aR_B(s)^2}\|f\|_{L^2(\mu)}^2.
 \label{vac:discardedform}
\end{equation}
Since $t=2s/x$ and $W_B=2bp_B/x$, this is precisely
\eqref{vac:heatbound}.

To justify the form restriction, observe that $P$ maps smooth functions to
smooth functions.  On each fixed compact graph it is bounded on $H^1(M)$:
differentiating its conditional integral in an exterior direction produces
the conditional derivative of $f$ and a covariance with
$2\nabla_y\log\psi$.  These coefficients are smooth and bounded, and
conditional Cauchy--Schwarz gives the required fixed-graph bound.  Derivatives
in the $B$ directions vanish for $Pf$.  Thus $Q$ is also bounded on
$H^1(M)$.  If $f\in H^1(M)\cap\operatorname{Ran}Q$, approximate $f$ in
$H^1$ by smooth $f_j$ and use $Qf_j\to f$.  This proves smooth density in
the restricted form domain.  Smooth density in $L^2$ and boundedness of $Q$
also prove that the restriction is densely defined in
$\operatorname{Ran}Q$.  It is closed because this subspace is closed in
$L^2(\mu)$ and $k$ is closed.  Inequality~\eqref{vac:discardedform} consequently
extends to the entire restricted form domain without any uniform derivative
bound on $\psi$.

For a fixed gauge transformation, the selected and exterior link coordinates
transform separately.  Both transformations preserve their Haar measures,
and $\psi$ is invariant.  Changing variables in
\eqref{vac:conditional} therefore shows that $P$ intertwines the gauge
action.  The forms and their restrictions are gauge invariant, proving the
physical-sector assertion.

Finally, if $bp_B=0$, the Hamiltonian splits as a free product Laplacian on
$B$ plus an exterior operator.  Its positive ground function is constant on
$B$.  The conditional measures are exactly product Haar, and
\eqref{vac:haar} gives~\eqref{vac:freebound}.
\end{proof}

\subsection{An explicit estimate at weak coupling}
\label{vac:explicit}

Let $D_N$ be the diameter of $\SU(N)$ in the specified Casimir metric.
The next heat-kernel estimate supplies explicit constants in
Theorem~\ref{vac:finiteblock}.

\begin{lemma}[Compact heat-kernel ratio]
\label{vac:heatlemma}
For every $s>0$,
\begin{equation}
 \kappa_N(s)\leq 2^n
       \exp\!\left(\frac{3D_N^2}{4s}\right),
 \qquad n=N^2-1.
 \label{vac:kappabound}
\end{equation}
\end{lemma}

\begin{proof}
The bi-invariant metric has nonnegative Ricci curvature.  For a positive
solution of $u_t=\Delta u$ on a compact $n$-dimensional manifold with
nonnegative Ricci curvature, the Li--Yau estimate~\cite{LiYau} is
\begin{equation}
 |\nabla\log u|^2-\partial_t\log u\leq\frac{n}{2t}.
 \label{vac:liyau}
\end{equation}
Here is its compact, boundaryless derivation.  Set $f=\log u$,
$w=|\nabla f|^2-f_t$, and $\mathcal F=tw$.  The heat equation, Bochner
identity, and $|\operatorname{Hess}f|^2\geq(\Delta f)^2/n$ imply
\begin{equation}
 (\Delta-\partial_t)\mathcal F
 \geq\frac{2\mathcal F^2}{nt}-\frac{\mathcal F}{t}
       -2\langle\nabla f,\nabla\mathcal F\rangle.
 \label{vac:bochnerheat}
\end{equation}
For smooth positive initial data, $\mathcal F=0$ at $t=0$.  A positive
space-time maximum exceeding $n/2$ contradicts~\eqref{vac:bochnerheat}
and the maximum principle.  This proves~\eqref{vac:liyau}.  For the heat
kernel, start after a positive time shift and then let that shift tend to
zero.  Compactness removes boundary and escape issues.

Along a constant-speed minimizing geodesic $\gamma$ between times
$0<t_1<t_2$, inequality~\eqref{vac:liyau} gives
\[
 \frac{\mathrm d}{\mathrm dt}f(\gamma(t),t)
 \geq-\frac{|\dot\gamma(t)|^2}{4}-\frac{n}{2t}.
\]
Integration yields
\begin{equation}
 u(z,t_1)\leq u(w,t_2)
       \left(\frac{t_2}{t_1}\right)^{n/2}
       \exp\!\left(\frac{\operatorname{dist}(z,w)^2}
                         {4(t_2-t_1)}\right).
 \label{vac:harnack}
\end{equation}
Apply this to $k_t(v,\cdot)$ between $s$ and $2s$ and integrate the later
point $w$ against normalized Haar measure.  Total heat-kernel mass one gives
\begin{equation}
 \max k_s\leq 2^{n/2}e^{D_N^2/(4s)}.
 \label{vac:heatupper}
\end{equation}
Apply~\eqref{vac:harnack} between $s/2$ and $s$ and integrate the earlier
point $z$ instead.  The result is
\begin{equation}
 \min k_s\geq2^{-n/2}e^{-D_N^2/(2s)}.
 \label{vac:heatlower}
\end{equation}
The two inequalities prove~\eqref{vac:kappabound}.  Normalizing Riemannian
volume to Haar probability multiplies the heat kernel by a constant and
does not change the logarithmic differential estimate; its mass in the
normalized measure is one, as used above.
\end{proof}

\begin{corollary}[Explicit finite-block lower bound]
\label{vac:exponential}
If $bp_B>0$, then
\begin{equation}
 D_B\geq\frac{xC_F}{2a}\,2^{-2mn}
                   e^{-C_B/x}I,
 \qquad C_B=4D_N\sqrt{3mbp_B}.
 \label{vac:explicitbound}
\end{equation}
\end{corollary}

\begin{proof}
Substituting~\eqref{vac:kappabound} into~\eqref{vac:heatbound} gives the
lower certificate
\[
 \frac{xC_F}{2a}\,2^{-2mn}
 \exp\!\left(-\frac{3mD_N^2}{2s}-\frac{8sbp_B}{x^2}\right).
\]
The positive exponent is minimized at
\[
 s_* =\frac{xD_N}{4}\sqrt{\frac{3m}{bp_B}},
 \qquad
 \frac{3mD_N^2}{2s_*}+\frac{8s_*bp_B}{x^2}
       =\frac{4D_N\sqrt{3mbp_B}}{x}.
\]
This proves~\eqref{vac:explicitbound}.  The case $bp_B=0$ is already covered
by~\eqref{vac:freebound} and does not require this optimization.
\end{proof}

The exponential dependence is a rigorous lower certificate, not an
asymptotic formula for the actual discarded spectrum as $x\downarrow0$.
In particular, the estimate is not uniform in $m$ or $N$.

\subsection{Conditional scale budget and the remaining limitation}
\label{vac:budget}

One can quantify the strength of this certificate without asserting a
renormalization theorem.  Suppose, as an additional hypothesis, that the same
fixed-block estimate with constants
$c_0=(C_F/2)2^{-2mn}$ and $C=C_B$ were available for every relevant exact
coarse operator.  Consider the illustrative scale and coupling sequences
\begin{equation}
 a_r=A L^{-r},\qquad x_r=(\sigma+qr)^{-1},\qquad r\geq1,
 \quad A>0,\ L>1,\ \sigma>0,\ q\geq0.
 \label{vac:modelscales}
\end{equation}
The reciprocal costs of the certified bounds
$d_r^{\mathrm{cert}}=c_0x_r e^{-C/x_r}/a_r$ then satisfy the exact identity
\begin{equation}
 \sum_{r\geq1}\frac1{d_r^{\mathrm{cert}}}
 =\frac{A}{c_0}e^{C\sigma}
       \sum_{r\geq1}(\sigma+qr)\rho^r,
 \qquad \rho=\frac{e^{Cq}}{L}.
 \label{vac:budgetseries}
\end{equation}
This series is finite exactly when $\rho<1$, in which case it equals
\begin{equation}
 \frac{A}{c_0}e^{C\sigma}
 \left(\frac{\sigma\rho}{1-\rho}
       +\frac{q\rho}{(1-\rho)^2}\right).
 \label{vac:budgetvalue}
\end{equation}
No actual Yang--Mills coupling trajectory, choice of scale factor, or
stability of the finite-block theorem under exact elimination is proved by
this calculation.  If $\rho\geq1$, this particular certificate does not
give a finite upper bound on the inverse-gap cost; the true discarded gaps
could still give a finite sum.

Theorem~\ref{vac:finiteblock} concerns one fixed finite block in the original
Hamiltonian~\eqref{vac:H}.  A whole elimination layer generally contains an
extensive number of links.  If that layer is used as $B$, the factor
$2^{-2mn}$ deteriorates exponentially in $m$, and $C_B$ grows proportionally
to $m$ when $p_B$ does.  Eliminating bounded blocks one after another instead
requires estimates for the successive exact coarse operators and control of
the number of elimination steps.  Neither follows from the theorem: the
generated operators have not been shown to retain the form~\eqref{vac:H},
and conditional bounds for separate blocks do not by themselves provide a
collective conditional Poincar\'e inequality in an interacting vacuum.
Thus the proved exterior-volume uniformity leaves open a whole-layer
estimate for the actual compact-link ground measure that persists under
exact coarse graining.
